\documentclass[11pt,a4paper]{article}
\usepackage[utf8]{inputenc}
\usepackage[T1]{fontenc}
\usepackage[french]{babel}
\usepackage{xcolor}
\usepackage{geometry}
\usepackage{amsmath}
\usepackage{fancybox}
\geometry{margin=2cm}

% Définition d'un environnement pour les algorithmes
\newcounter{algocounter}
\newenvironment{algorithm}[1]{
  \refstepcounter{algocounter}
  \begin{center}
  \shadowbox{
    \begin{minipage}{0.9\textwidth}
      \textbf{Algorithme \thealgocounter : #1}
      \vspace{0.2cm}
      \hrule
      \vspace{0.3cm}
}{
    \end{minipage}
  }
  \end{center}
  \vspace{0.3cm}
}

% Commandes pour le pseudo-code
\newcommand{\KW}[1]{\textbf{#1}}
\newcommand{\COMMENT}[1]{\textcolor{gray}{\textit{// #1}}}
\newcommand{\INDENT}{\quad}
\newcommand{\FUNCTION}[2]{\KW{Fonction} #1(#2)}
\newcommand{\PROCEDURE}[2]{\KW{Procédure} #1(#2)}
\newcommand{\IF}[1]{\KW{Si} #1 \KW{alors}}
\newcommand{\ELSE}{\KW{sinon}}
\newcommand{\ELSEIF}[1]{\KW{sinon si} #1 \KW{alors}}
\newcommand{\ENDIF}{\KW{fin si}}
\newcommand{\FOR}[1]{\KW{Pour} #1 \KW{faire}}
\newcommand{\ENDFOR}{\KW{fin pour}}
\newcommand{\WHILE}[1]{\KW{Tant que} #1 \KW{faire}}
\newcommand{\ENDWHILE}{\KW{fin tant que}}
\newcommand{\RETURN}[1]{\KW{retourner} #1}
\newcommand{\STATE}[1]{#1}

\title{Exercices d'Algorithmique}
\author{}
\date{\today}

\begin{document}

\maketitle

\section{Exercice 1 : Calcul de la Factorielle}

\textbf{Énoncé :} Écrire un algorithme qui calcule la factorielle d'un nombre entier positif n.

\begin{algorithm}[H]
\caption{Calcul de la factorielle (version itérative)}
\begin{algorithmic}[1]
\Function{Factorielle}{n}
    \State résultat $\leftarrow$ 1
    \For{i de 1 à n}
        \State résultat $\leftarrow$ résultat $\times$ i
    \EndFor
    \State \Return résultat
\EndFunction
\end{algorithmic}
\end{algorithm}

\section{Exercice 2 : Vérification de Nombre Premier}

\textbf{Énoncé :} Écrire un algorithme qui détermine si un nombre est premier.

\begin{algorithm}[H]
\caption{Test de primalité}
\begin{algorithmic}[1]
\Function{EstPremier}{n}
    \If{n $\leq$ 1}
        \State \Return faux
    \EndIf
    \If{n = 2}
        \State \Return vrai
    \EndIf
    \For{i de 2 à $\sqrt{n}$}
        \If{n mod i = 0}
            \State \Return faux
        \EndIf
    \EndFor
    \State \Return vrai
\EndFunction
\end{algorithmic}
\end{algorithm}

\section{Exercice 3 : Algorithme de Tri à Bulles}

\textbf{Énoncé :} Implémenter l'algorithme de tri à bulles.

\begin{algorithm}[H]
\caption{Tri à bulles}
\begin{algorithmic}[1]
\Procedure{TriBulles}{T, n}
    \For{i de 1 à n-1}
        \For{j de 1 à n-i}
            \If{T[j] > T[j+1]}
                \State échanger T[j] et T[j+1]
            \EndIf
        \EndFor
    \EndFor
\EndProcedure
\end{algorithmic}
\end{algorithm}

\section{À Compléter}

\subsection{Exercice 4 : Calcul du PGCD}
\textbf{Énoncé :} Écrire un algorithme pour calculer le PGCD de deux nombres (algorithme d'Euclide).

\begin{algorithm}[H]
\caption{Votre solution ici}
\begin{algorithmic}[1]
\Function{PGCD}{a, b}
    \State \Comment{À compléter}
\EndFunction
\end{algorithmic}
\end{algorithm}

\subsection{Exercice 5 : Recherche dans un Tableau}
\textbf{Énoncé :} Écrire un algorithme de recherche linéaire dans un tableau.

\begin{algorithm}[H]
\caption{Votre solution ici}
\begin{algorithmic}[1]
\Function{RechercheLineaire}{T, n, valeur}
    \State \Comment{À compléter}
\EndFunction
\end{algorithmic}
\end{algorithm}

\end{document}